\labsection{1}{Классификация с помощью персептрона}

\labpart{Цель работы}
Изучить модель персептрона и правило его обучения, исследовать возможности и ограничения однослойного
персептрона на задачах бинарной и многоклассовой классификации.

\labpart{Теоретические сведения}
Подразделы~\ref{sec:ml}.3 и~\ref{sec:perceptron}.1--\ref{sec:perceptron}.3 теоретической части: модель
нейрона с пороговой функцией активации, правило обучения (\ref{eq:perceptron-rule}), теорема
сходимости, проблема XOR, кодирование классов выходами нескольких нейронов.

\labpart{Постановка задачи}
Исходные данные заданы вариантом (таблица~\ref{tab:l1-var}): четыре точки плоскости и их классы $T$ для
задания~1, логическая функция для задания~2, центры и разброс кластеров для задания~3.
\begin{steps}
\item Построить персептрон, разделяющий на два класса точки варианта.
\item На векторах $(0; 0)$, $(0; 1)$, $(1; 0)$, $(1; 1)$ с классами, заданными логической функцией
варианта (XOR, исключающее ИЛИ~--- $0, 1, 1, 0$; XNOR, эквивалентность~--- $1, 0, 0, 1$),
продемонстрировать невозможность линейного разделения.
\item Построить персептрон из двух нейронов, разбивающий входные векторы на четыре класса: четыре
группы по 25 точек с нормальным разбросом $\sigma$ вокруг центров варианта. Коды классов $00$, $01$, $10$,
$11$ назначить так, чтобы каждый разряд кода был линейно разделим.
\end{steps}

\begin{table}[!htb]
\caption{Варианты заданий к лабораторной работе №1}
\label{tab:l1-var}
\centering\footnotesize\setlength{\tabcolsep}{3.5pt}
\begin{tabular}{clcclc}
\toprule
№ & Задание 1: точки $(x_1; x_2)$ & Классы $T$ & Задание 2 & Задание 3: центры & $\sigma$ \\
\midrule
1 & $(0,1; -0,1)$, $(-0,5; 0,6)$, $(-0,3; -0,9)$, $(0,9; -0,9)$ & 0 0 1 1 & XOR & $(\pm 1,5; \pm 1,5)$ & 0,3 \\
2 & $(0,6; 0,9)$, $(1; 0)$, $(-0,7; -1)$, $(0; -0,8)$ & 0 1 0 1 & XNOR & $(\pm 2; 0)$, $(0; \pm 2)$ & 0,3 \\
3 & $(-0,5; -0,5)$, $(-0,5; 0,5)$, $(0,3; -0,5)$, $(-0,1; 1)$ & 0 0 1 1 & XOR & $(\pm 2; \pm 2)$ & 0,3 \\
4 & $(-0,8; 0,9)$, $(-0,8; -0,7)$, $(0,6; -0,4)$, $(-0,2; -0,2)$ & 1 0 0 1 & XNOR & $(\pm 2,5; 0)$, $(0; \pm 1,5)$ & 0,3 \\
5 & $(0,2; 0,2)$, $(0,9; -0,8)$, $(-0,6; 0,3)$, $(-0,3; -0,7)$ & 1 1 0 0 & XOR & $(\pm 2; \pm 1,5)$ & 0,3 \\
6 & $(0,4; 0,8)$, $(0,9; 0,2)$, $(-0,3; 0,4)$, $(0,9; 1)$ & 0 1 1 0 & XNOR & $(\pm 1,5; 0)$, $(0; \pm 2,5)$ & 0,3 \\
7 & $(-0,1; 0,6)$, $(-0,4; -0,9)$, $(0,5; 0,6)$, $(0,7; -0,1)$ & 1 0 1 0 & XOR & $(\pm 1,5; \pm 2)$ & 0,3 \\
8 & $(1; -1)$, $(-0,5; 0,9)$, $(0,9; -0,4)$, $(1; 0,6)$ & 0 0 1 1 & XNOR & $(\pm 2,5; 0)$, $(0; \pm 2)$ & 0,4 \\
9 & $(0,6; 0,1)$, $(-0,2; -1)$, $(-0,3; -0,3)$, $(-0,8; -0,8)$ & 0 1 0 1 & XOR & $(\pm 2,5; \pm 1,5)$ & 0,4 \\
10 & $(0,8; 0,9)$, $(0,1; 0,3)$, $(0,9; 0)$, $(-0,5; 0,9)$ & 1 0 0 1 & XNOR & $(\pm 3; 0)$, $(0; \pm 2)$ & 0,4 \\
11 & $(0,8; 1)$, $(-0,4; 0,7)$, $(1; -0,4)$, $(-1; -0,7)$ & 1 1 0 0 & XOR & $(\pm 1,5; \pm 2,5)$ & 0,4 \\
12 & $(-0,2; -0,6)$, $(-0,9; -0,9)$, $(0,7; -1)$, $(0,3; 0,2)$ & 0 1 1 0 & XNOR & $(\pm 2,5; 0)$, $(0; \pm 2,5)$ & 0,4 \\
13 & $(0,5; -0,1)$, $(-0,9; -0,9)$, $(1; 0,4)$, $(0; -0,7)$ & 1 0 1 0 & XOR & $(\pm 2,5; \pm 2,5)$ & 0,5 \\
14 & $(0,2; -1)$, $(0,4; 0,9)$, $(-0,2; 0,7)$, $(-0,5; 0,1)$ & 0 0 1 1 & XNOR & $(\pm 2; 0)$, $(0; \pm 3)$ & 0,4 \\
\bottomrule
\end{tabular}
\end{table}

Код обучения на NumPy приведён в подразделе~\ref{sec:py-perceptron}; точки классов выводятся функцией
\code{plt.scatter}, разделяющая прямая $x_2 = -(w_1x_1 + b)/w_2$~--- функцией \code{plt.plot}.

\labpart{Требования к интерфейсу и анимации}
\begin{itemize}
\item \emph{Графическое поле.} Двумерная координатная сетка в диапазоне $[-1; 1]$ по обеим осям (в задании~3~--- $[-4; 4]$); точки классов показаны разными маркерами и цветами, поверх них~--- текущая разделяющая прямая. Рядом выводятся график числа ошибок по эпохам и текущие значения $w_1$, $w_2$, $b$.
\item \emph{Интерактивность.} Элементы управления: коэффициент обучения $\eta$, максимальное число эпох, функция активации (пороговая $0/1$ или знаковая $\pm 1$), кнопки <<Старт>>, <<Пауза>>, <<Шаг>>, <<Сброс>>. Новые точки добавляются на график щелчком мыши с выбором класса.
\item \emph{Мультипликация.} Разделяющая прямая меняет наклон и смещение непосредственно в процессе обучения~--- после каждой коррекции весов, а не один раз за эпоху. Ползунок <<Задержка (мс)>> замедляет отрисовку шагов.
\item \emph{Подсветка событий.} Точка, на которой персептрон ошибся на текущем шаге, кратковременно выделяется миганием или увеличением маркера. В задании~3 плоскость закрашивается по коду $(y_1, y_2)$, и раскраска обновляется вместе с прямыми.
\end{itemize}

\labpart{Порядок выполнения работы}
\begin{steps}
\item Реализовать персептрон на NumPy (подраздел~\ref{sec:py-perceptron}).
\item Выполнить задание~1 с перерисовкой разделяющей прямой после каждой эпохи; зафиксировать итоговые
веса и число эпох.
\item Выполнить задание~2 с ограничением в 50 эпох, построить график числа ошибок по эпохам и объяснить
результат.
\item Выполнить задание~3: обучить оба нейрона, построить две разделяющие прямые и карту решений~---
раскраску плоскости по коду $(y_1, y_2)$; вычислить долю верно классифицированных точек.
\end{steps}

\labpart{Пример выполнения}
Результаты для варианта~3 приведены на рисунке~\ref{fig:l1}. В задании~1 персептрон с нулевыми
начальными весами и $\eta = 1$ разделил классы за 3 эпохи (в третьей эпохе ошибок нет):
$\mathbf{w} = (2; 0,5)$, $b = 0$. В задании~2 после двух эпох персептрон ошибается на всех четырёх
точках XOR: коррекции компенсируют друг друга, к концу каждой эпохи веса возвращаются к значениям
$\mathbf{w} = (-1; 0)$, $b = 0$, и обучение не сходится.
В задании~3 два нейрона разделили 100 точек на четыре класса за 2 эпохи без ошибок: каждый нейрон
отвечает за один разряд кода и строит одну из двух прямых.

\begin{figure}[!htb]
\centering
\includegraphics[width=\textwidth]{\figdir/lab1_example.png}
\caption{Персептрон: разделение двух классов, число ошибок на задаче XOR, карта решений персептрона из
двух нейронов}
\label{fig:l1}
\end{figure}

\labpart{Контрольные вопросы}
\begin{questions}
\item Опишите модель персептрона Розенблатта. Какую функцию активации он использует?
\item Какое геометрическое место точек задаёт уравнение $\mathbf{w}^{\mathsf T}\mathbf{x} + b = 0$? Как
расположен вектор весов относительно разделяющей прямой и за что отвечает смещение?
\item Запишите правило обучения персептрона. Как изменяются веса при ошибке каждого знака?
\item Сформулируйте теорему сходимости персептрона. Единственна ли найденная граница?
\item Какие множества называются линейно разделимыми? Докажите, что функция XOR линейно неразделима.
Как решить эту задачу нейронной сетью?
\item Сколько классов может различить слой из $m$ персептронов? Какие требования предъявляются к
кодированию классов?
\end{questions}
